\section{FinTrace 智能体架构与执行机制}
\label{sec:architecture}

\subsection{总体技术架构}

面向金融长上下文问答中多轮语义继承、动态工具调用、复杂任务拆解以及事实可追溯等需求，FinTrace构建了一套以确定性控制流为骨架、以大语言模型为语义决策核心、以专业金融工具为事实执行单元的证据驱动型智能体架构。系统并未采用多个自由智能体相互协商的开放式工作模式，而是将大语言模型的决策空间约束在预先定义的能力边界和状态转换体系之内，使复杂金融问题的处理过程能够被记录、校验和复现。

从整体技术路径看，FinTrace采用“LangGraph确定性编排骨架+Direct Fast Path+Evidence-guided Bounded Investigation”的混合架构。大语言模型主要承担复杂语义理解、任务识别以及下一步调查动作规划，财务计算、股权路径搜索、事件筛选、文档检索、时间过滤和停止条件等确定性任务则由程序及专业工具完成。工具产生的结构化结果经过校验后统一转换为Evidence，并由Evidence Ledger管理，最终回答只能基于已经验证的Evidence形成。由此，系统将大语言模型的开放推理能力嵌入可验证的执行框架之中，在保留复杂金融研究能力的同时降低事实性幻觉、错误工具调用和无依据推断的风险。

当前系统以LangGraph \texttt{StateGraph}作为在线Agent的统一控制骨架。CLI、FastAPI 与Next.js Web前端均接入同一套执行工作流，因此本地调试、接口调用、前端演示和后续评测使用相同的状态定义和运行逻辑。工具、数据源与大语言模型之间通过统一Schema建立交互契约，只要输入输出协议保持一致，底层检索器、模型或数据索引均可独立更新，而无需重构整个Agent控制流程。FinTrace 当前总体架构的主要层次及职责如表~\ref{tab:architecture_1} 所示。

\begin{table}[htbp]
  \centering
  \caption{FinTrace总体架构主要层次及职责}
  \label{tab:architecture_1}
  \begin{tabular}{C{.23\textwidth}C{.67\textwidth}}
    \fintracetoprule
    \fintracetablehead{架构层级} & \fintracetablehead{主要职责} \\
    \fintracemidrule
    交互接入层 & 通过CLI、FastAPI与Next.js Web前端接收用户请求，并展示回答、工具执行状态、Evidence与运行Trace \\
    Agent编排层 & 完成会话恢复、请求解析、可回答性判断、任务路由、动作校验、有界调查、证据审查和回答生成 \\
    专业工具层 & 提供财务分析、股权分析、事件时间线、研报观点分析及文档检索五类独立能力 \\
    数据与索引层 & 管理标准化金融数据、股东快照、公告事件、研报观点、Document、Chunk，以及SQLite、FAISS和FTS5/BM25等索引 \\
    可信控制层 & 通过Evidence Ledger、Knowledge Cutoff、Action Validation、Trace和Evaluation约束完整执行过程 \\
    \fintracebottomrule
  \end{tabular}
\end{table}

这种分层结构使FinTrace不再依赖“大模型直接给出最终答案”的单一生成范式，而是将金融问答转换为一系列具有明确输入、动作、证据和终止条件的可执行任务。交互层负责接收问题，Agent编排层负责判断下一步应该做什么，专业工具负责回答事实层面的“数据是什么”，Evidence体系负责判断“这些结果是否足以支持最终结论”。不同层次之间保持职责分离，从而为后续长周期记忆、动态路由和防幻觉机制提供统一技术基础。

\subsection{从请求解析到有界调查}

FinTrace在接收到用户问题后并不立即调用大语言模型生成结论，而是首先恢复当前Session的历史状态，并将自然语言问题转换为可执行的结构化请求。当前在线工作流依次完成会话恢复、请求解析、可回答性判断、任务路由、动作校验、工具执行、工具结果验证、证据合并与证据充分性审查，并根据当前Evidence状态决定继续调查或进入最终回答阶段。

请求解析阶段主要解决“用户正在问什么”这一问题。系统综合当前输入与历史上下文识别当前主体、人物、比较对象、时间锚点、时间范围、财务报告期、指标、任务类型、能力缺口及潜在歧义。本轮明确出现的实体与任务优先于历史状态，当用户仅通过简称、股票代码或代词延续上一轮查询时，系统再结合\texttt{current\_context}判断是否存在可以可靠继承的上下文信息。

时间信息在请求解析中被单独处理。系统区分明确日期、今日、最新、近期以及未指定时间等不同语义，并将统一知识截止日作为数据访问边界。知识截止日由程序注入，Planner无权自行修改，因此可以防止模型在规划过程中使用超出比赛数据范围的未来信息。对于财务任务，报告期还需进一步经过确定性期间解析器转换，以避免大语言模型自行猜测年份、混合不同累计口径或将知识截止日错误地视为财务报告期。

完成解析后，系统执行Pre-Answerability判断。当前实现并非简单区分“可以回答”与“不可以回答”，而是进一步划分为四类状态。

\begin{table}[htbp]
  \centering
  \caption{Pre-Answerability状态及处理逻辑}
  \label{tab:architecture_2}
  \begin{tabular}{C{.18\textwidth}C{.27\textwidth}C{.45\textwidth}}
    \fintracetoprule
    \fintracetablehead{状态} & \fintracetablehead{判定含义} & \fintracetablehead{后续处理} \\
    \fintracemidrule
    routeable & 系统具备相应能力，且执行所需关键参数完整 & 进入正常路由 \\
    routeable\_with\_gaps & 请求包含可执行内容，同时存在非阻断能力缺口 & 继续调查可支持部分，并在最终回答中披露未覆盖内容 \\
    clarification\_required & 主体、任务或比较对象等关键条件无法唯一确定 & 暂停工具调用并请求用户补充信息 \\
    unsupported & 当前系统不存在对应数据或工具能力 & 明确说明能力边界，不调用无关工具替代 \\
    \fintracebottomrule
  \end{tabular}
\end{table}

这种设计能够避免在能力缺失时将问题错误映射到“最相似”的工具。例如，当一个问题同时包含财务分析和实时股价信息时，系统可以继续完成财务部分，并将实时行情记录为capability gap，而不是因为缺少实时行情接口而放弃整条任务，也不会将行情问题错误路由至财务工具。

在完成可回答性判断后，FinTrace根据任务复杂程度选择Direct Fast Path或Bounded Investigation。对于实体、时间、指标和操作均能够唯一解析的简单事实请求，Direct Gate可以直接形成确定性的工具Action。例如，用户明确询问某公司某一报告期的营业收入时，只要主体、指标与报告期均完成解析，系统即可直接调用对应的财务指标查询操作，从而避免不必要的Planner调用。

Direct Fast Path并不意味着跳过可信控制。即使是简单的确定性调用，其工具结果仍需经过结果合法性检查、Evidence转换、Evidence Ledger合并及证据充分性判断后，才能进入最终回答环节。

当问题涉及多步骤调查、多个证据来源或原因追溯时，系统进入Evidence-guided Bounded Investigation。Planner在每一轮仅产生一个AgentAction，并明确本轮需要调用的工具、具体Operation、执行参数、调用目的和预期获得的Evidence。工具执行结束后，系统重新评估现有证据是否已经足以覆盖用户问题。若仍存在可以通过当前能力补足的核心信息缺口，则继续规划下一项动作；若证据已经充分，则终止调查并生成回答；若已经无法获得关键证据，则进入部分回答或insufficient\_evidence状态。

为了防止Agent在复杂问题上形成无限循环，调查过程设置了明确的执行预算，包括最大调查步骤数、最大工具调用总数、重复调用限制、文档检索次数以及连续无新增Evidence的终止条件。因此，FinTrace所采用的自主调查是一种具有明确边界和停止条件的有限搜索过程，而不是无限扩展的ReAct式推理。

\subsection{长周期记忆、动态路由与受控自纠错}

长周期金融问答的主要难点并不只是对话内容变长，而是随着轮次增加，历史信息中真正与当前任务相关的部分所占比例不断降低。如果每一轮都将全部历史内容重新输入大语言模型，不仅会显著增加上下文成本，也容易使已失效的实体、时间和旧任务持续干扰当前判断。FinTrace因此采用分层记忆方式维护Session状态，将近期原始消息、结构化上下文、滚动摘要与已验证事实分别承担不同职责。

近期原始消息主要保留距离当前轮最近的完整用户输入和Agent回答，用于解析局部指代、简称和连续追问。当前实现的常规消息窗口上限为12条。当达到固定轮次、消息窗口超过设定规模或累计字符数超过阈值时，系统启动memory\_summarizer，将即将离开窗口的历史信息与既有摘要进行合并。摘要成功后保留最近8条原始消息；如果摘要调用失败，则继续保留最近12条并记录警告，不影响已经完成的当前回答。

结构化CurrentContext保存当前公司、人物、比较对象、时间、指标和任务等信息，用于维持多轮查询中的明确状态。相比仅依赖自然语言历史，结构化状态能够更稳定地处理“那另一家公司呢”“它的评级呢”等省略式表达。本轮出现的新实体和明确任务优先于历史上下文，只有在当前信息不足且历史状态能够唯一解释时，系统才进行语义继承。

Conversation Summary负责压缩较早的对话内容，使系统能够在不持续回填全部历史原文的情况下保留较长期语境。该摘要用于恢复上下文关系，不被直接视为事实证据。

Verified Findings则属于更严格的已验证事实记忆。其内容不是从最终自然语言回答反向提取，而是直接来自Evidence Ledger。写入该区域的信息必须具有有效Evidence ID，同时保留对应公司、主题、来源工具和来源轮次。大语言模型的猜测、失败工具结果以及没有证据编号支持的结论均不得进入Verified Findings。在后续对话中，这些已验证信息可以作为Planner恢复研究语境和选择核验路径的依据，但如果相关事实需要再次成为当前回答中的核心结论，系统仍要求本轮工具重新取得Evidence，从而避免“记忆中的旧结论”替代当前证据。

面向赛题提出的0.5M Tokens长上下文目标，FinTrace当前采取的技术思路并不是不断扩大单轮Prompt，而是通过消息窗口、结构化状态、滚动摘要和验证事实进行分层压缩与按需恢复。现阶段系统已经实现上述分层记忆和会话持久化机制，但50K—500K Tokens范围的正式冻结压力测试尚未完成，因此当前白皮书仅将其描述为面向超长上下文目标的已实现架构基础，而不直接表述为0.5M Tokens技术指标已经达到。后续将通过统一压力测试进一步验证关键事实召回率、任务连续性和回答准确率。

在工具路由方面，FinTrace同样不采用简单的关键词映射方式，而是根据当前任务、主体状态、时间约束、已有Evidence和Capability Gap共同确定下一步动作。当前系统提供财务分析、股权分析、事件时间线、研报观点分析和文档检索五类工具。财务工具负责指标查询、比较及风险扫描；股权工具负责股东快照、变化比较和有限持股路径；事件工具负责事件节点查询及事件聚合；研报观点工具负责机构评级、盈利预测、风险提示与分析师判断；文档检索则主要承担公告、研报及其他金融文本的原文证据补充。

工具之间采用分层协同，而不是互相替代。例如，用户询问“证券机构如何评价某家公司”时，系统优先通过research\_analysis获取带有机构、日期、观点类型和来源的结构化结果；只有当用户进一步追问“为什么得出这一判断”或“原文具体如何描述”时，才调用document\_search获取对应研报Chunk。对于事件问题，系统先通过event\_timeline定位结构化事件，再根据用户对金额、原因、整改细节等正文信息的需求调用文档检索。该机制减少了大模型直接面对大规模无关文本的需要，也使不同类型证据的语义边界更加明确。

为了控制Planner的错误动作，FinTrace在工具执行前设置了Action Validator。该模块检查工具名称、Operation、Schema参数、目标实体、报告期、知识截止日、重复调用及调查预算。当Planner输出非法动作时，系统只允许进行有限修复或重新规划，不进入无限反思循环。工具执行后，系统继续区分成功、可重试失败、不可重试失败、DATA\_NOT\_AVAILABLE、insufficient\_data与not\_applicable等状态，从而避免将“当前数据没有命中”误写成“现实世界中不存在”。

由此，FinTrace的自纠错机制表现为一个受控制的工程闭环。系统首先识别动作或工具结果中的异常，再根据错误类型决定是否允许修复或重新规划，修复后的Action重新进入合法性检查，超出预算或无法恢复的情况则明确终止。整个过程中产生的原始Action、校验错误、修复结果、工具返回、证据缺口和终止原因均记录在Trace中，为后续工具调用Precision、自纠错成功率以及失败案例分析提供可直接复用的数据基础。

\subsection{Evidence约束与可追溯回答机制}

FinTrace将Evidence作为连接金融工具与最终自然语言回答的核心中间状态。防幻觉并不是通过在Prompt中简单要求模型“不要编造”来实现，而是通过对“哪些内容能够进入最终回答”建立明确的程序约束。

每个经过工具结果验证的有效结果都会被转换为具有唯一evidence\_id的事实单元，并写入Evidence Ledger。Evidence中保留来源工具、Operation、目标对象、对应时间、原始数据或文本定位、必要的计算输入和公式、质量标记以及适用限制。对于文档类结果，Evidence可以进一步追溯至chunk\_id、document\_id和原始材料位置；对于财务、股权、事件及研报观点结果，则能够回到对应结构化数据源和工具执行记录。

Document与Chunk的设计进一步保证了文本证据的稳定性。比赛语料首先被转换为规范化Document，再按照标题、段落、列表和表格等语义边界形成冻结Chunk。每个Chunk均具有稳定的ID、内容哈希、语料版本和来源位置。在线检索先按照证券主体、文档类型、报告期、发布日期和知识截止日进行结构化过滤，再融合FTS5/BM25词法召回与向量相似度检索。由此，系统不仅能够找到与问题语义接近的文本，还能够将最终引用定位到确定的源文档和文本范围。

Evidence Ledger中的内容在进入Final Answer之前还需要接受Evidence Review。Reviewer重点判断当前Evidence是否已经覆盖用户问题中的核心子任务，以及不同来源之间是否存在明显冲突。当核心问题已经具备充分且一致的证据时，系统进入回答生成；当部分内容已经得到支持、部分内容仍存在非阻断缺口时，系统允许形成部分回答，并明确说明尚未覆盖的部分；当关键结论缺少证据且剩余调查预算内也无法进一步补足时，则返回证据不足状态。

最终回答同时保存used\_evidence\_ids，用于标记真正支撑本轮结论的Evidence。因此，从最终文本可以沿着Evidence ID回溯至工具调用结果，并进一步定位至结构化数据记录、Document或Chunk，从而形成完整的证据链。

在最终表达阶段，系统进一步区分客观事实、系统推论、证券机构观点、风险信号和能力限制。客观事实必须能够直接回指Evidence；系统推论需要说明依据及不确定性；研报观点必须保留机构或分析师归属；财务风险信号只能解释为异常或预警，不能直接升级为财务造假结论；有限持股路径也不能表述为已覆盖全部工商信息的完整控制关系。对于当前没有实时行情、用户账户或其他必要数据支持的问题，系统明确披露能力边界，而不是利用相邻数据或语言常识生成替代性结论。

如果最终回答模型调用失败，系统进入Structured Error状态并返回明确错误信息，不使用固定文本模板模拟已经完成的智能分析。结合前述请求解析、Pre-Answerability、Action Validation、Tool Result Validation和Evidence Review，FinTrace最终形成了贯穿完整链路的可信控制体系，使金融智能问答从传统的一次性文本生成过程转变为任务可解析、动作可约束、工具可验证、证据可回溯、错误可修复和能力边界可解释的Agentic AI执行过程。

\begin{figure}[htbp]
  \centering
  \includegraphics[width=.92\textwidth]{word-ch2-flow.png}
  \caption{FinTrace 智能体执行流程}
  \label{fig:agent_workflow}
\end{figure}
